\chapter{Memoria de frontera: descomposición par--impar y \(3\) realizaciones hidrodinámica, gauge y cosmológica}
\label{chap:memoria-tres-lecturas-rev6}
\index{memoria de frontera}\index{hidrodinámica}
\index{Yang--Mills}\index{cosmología}

La unidad entre hidrodinámica, teoría gauge y cosmología no se obtiene
comparando tres ecuaciones una vez escritas.  Aparece antes, en la
descomposición de un mismo transporte en una parte par, que conserva el coste
de la ida y el retorno, y una parte impar, que conserva su orientación.  El
objeto común se construye en el sistema temporal dodecafásico \(\mathcal R_{12}\); sus representaciones
posteriores no se identifican entre sí.

\section{De los canales angulares al transporte dodecafásico}

Sean
\[
 q_\pm=\exp\!\left[-\frac{\pi}{180}(A\pm C^\ast)\right],
 \qquad
 S_{90}(q)=\frac{q^{90}}{1-q^{270}},
 \qquad
 S_{120}(q)=\frac{q^{120}}{1-q^{360}},
\]
y definamos
\[
 \kappa
 =12\bigl(S_{90}(q_-)-S_{90}(q_+)\bigr)
 -\bigl(S_{120}(q_-)-S_{120}(q_+)\bigr).
\]
Esta fórmula conserva la genealogía
\[
 \alpha_{\rm HMT}\longrightarrow A
 \longrightarrow C^\ast\longrightarrow q_\pm
 \longrightarrow\kappa.
\]
En particular, \(\kappa\) no es una viscosidad, una tensión ni una masa.  Es
el coeficiente adimensional de transferencia que esas magnitudes podrán leer
después.

Para
\[
 \phi_m=(30m+20)^\circ,\qquad
 p_m=12\cos(3\phi_m)-\cos(4\phi_m),
 \qquad 0\leq m<12,
\]
se ponen
\[
 w_m=\frac{1+\kappa p_m}{12}.
\]
La ortogonalidad de los caracteres de
\(\mathbb Z/12\mathbb Z\) da
\(\sum_m p_m=0\), luego \(\sum_mw_m=1\).  Además,
\(|p_m|\leq13\); para el valor construido de \(\kappa<1/13\), todos los
pesos son positivos.  Sobre
\(\mathcal H_{12}=\ell^2(\mathbb Z/12\mathbb Z)\) definimos
\[
 (U_{12}f)(n)=\sum_{m=0}^{11}w_mf(n-m).
\]

\begin{theorem}[Reserva dodecafásica exacta]
\label{thm:reserva-dodecafasica-rev6}
Sean
\[
 P_{12}=\frac{U_{12}+U_{12}^{\ast}}2,\qquad
 J_{12}=\frac{U_{12}-U_{12}^{\ast}}2,\qquad
 L_{12}=I-P_{12}.
\]
Entonces \(P_{12}\) es autoadjunto, \(J_{12}\) es
anti-autoadjunto y, para todo \(f\perp\mathbf1\),
\[
 \boxed{\;
 \langle f,L_{12}f\rangle
 \geq g_{12}\|f\|_2^2,\qquad
 g_{12}=1-3\kappa>0.
 \;}
\]
La constante es óptima y se alcanza en las frecuencias \(3\) y \(9\).
\end{theorem}

\begin{proof}
En la base de Fourier
\(e_k(n)=\exp(2\pi ikn/12)\), la convolución es diagonal:
\[
 U_{12}e_k=\lambda_k^{(U)}e_k,\qquad
 \lambda_k^{(U)}=\sum_{m=0}^{11}w_m e^{-2\pi ikm/12}.
\]
El modo constante da \(\lambda_0^{(U)}=1\).  Los cosenos de frecuencias tres y
cuatro sólo sobreviven en \(k=3,9\) y \(k=4,8\), respectivamente; los demás
modos se anulan por ortogonalidad.  El mayor valor de
\(\operatorname{Re}\lambda_k^{(U)}\), con \(k\neq0\), es \(3\kappa\).  Por tanto,
\[
 \langle f,L_{12}f\rangle
 =\sum_{k\neq0}(1-\operatorname{Re}\lambda_k^{(U)})|\widehat f(k)|^2
 \geq(1-3\kappa)\|f\|_2^2.
\]
\end{proof}

La descomposición posee un significado preciso.  \(P_{12}\) conserva aquello
que una ruta comparte con su inversa; \(J_{12}\) conserva lo que cambia de
signo al invertir la orientación.  El primero puede suministrar una forma
positiva.  El segundo no: para \(f\) real,
\(\langle f,J_{12}f\rangle\) es puramente imaginario.  A cambio,
\(J_{12}\) retiene la circulación que la proyección par perdería.

\section{Memoria distribuida y generador local}

La misma dualidad aparece en el núcleo
\[
 K_q(m,n)=q^{|m-n|},\qquad 0<q<1.
\]
Su matriz parece no local porque correlaciona niveles arbitrariamente
separados.  Sin embargo, sobre una cadena finita \(0,\ldots,N\), su inversa
es tridiagonal:
\[
 K_q^{-1}
 =\frac1{1-q^2}
 \begin{pmatrix}
 1&-q&0&\cdots&0\\
 -q&1+q^2&-q&\ddots&\vdots\\
 0&-q&1+q^2&\ddots&0\\
 \vdots&\ddots&\ddots&1+q^2&-q\\
 0&\cdots&0&-q&1
 \end{pmatrix}.
\]

\begin{proposition}[Green no local, acción local]
\label{prop:green-local-memoria-rev6}
Para \(0<q<1\), \(K_q\) es definida positiva y la matriz anterior es su
inversa.  En consecuencia, una memoria distribuida exponencialmente es el
Green de una acción de primeros vecinos.
\end{proposition}

\begin{proof}
Multiplicar la matriz tridiagonal por \(K_q\) cancela, en cada fila interior,
los términos
\[
 -q\,q^{|m-1-n|}
 +(1+q^2)q^{|m-n|}
 -q\,q^{|m+1-n|}
\]
cuando \(m\neq n\), y deja \(1-q^2\) cuando \(m=n\).  Las filas extremas
obedecen la misma identidad con un solo vecino.  La inversa es positiva;
por tanto \(K_q\) también lo es.
\end{proof}

Este resultado explica una aparente oposición que recorre Mecánica
Dimensional: la historia es extensa porque conserva correlaciones entre
profundidades, mientras el generador que la actualiza puede ser local.
``Local'' y ``con memoria'' no son alternativas.

\section{Representaciones locales interior y reflejada: circulación, holonomía, contracción y positividad de transferencia}

El capítulo de la hoja común ha demostrado una sola vez las dos
representaciones locales del borde. En la lectura interior, el sector impar
porta circulación y el sector par produce contracción y coercividad; en la
lectura reflejada, el primero porta holonomía y el segundo positividad de
transferencia y brecha finita. La identidad de Stokes y la coincidencia del
espectro no nulo de \(d^\ast d\) y \(dd^\ast\) son propiedad del objeto
común, no dos pruebas paralelas. Remitimos a
\cref{thm:operador-comun-frontera-ns-ym,thm:contraccion-flujo-finito,%
thm:coercividad-gap-gauge-finito,thm:hoja-comun-frontera}; aquí añadimos
únicamente la tercera lectura, la global.

\section{Lectura global: expansión y curvatura}

La tercera lectura no identifica el universo con una celda.  Toma la misma
separación entre interior y frontera como ley global de evolución.  La
identidad
\[
 1000=729+271
\]
recuerda que el crecimiento interior y la publicación de frontera completan
una unidad normalizada.  En una geometría homogénea e isótropa, la parte par
de un transporte puede representarse por una presión efectiva
\[
 p_{\rm eff}=p-3\zeta H,
\]
que es la forma tensorial estándar de una viscosidad de bulto.  La ecuación
no define \(\zeta\) desde \(H\); declara el tipo de la lectura cosmológica:
la memoria simétrica se publica como resistencia a la expansión.

La geometría de de Sitter muestra, sin ambigüedad, qué permanece común entre
tres foliaciones distintas.  Para \(H>0\),
\[
\begin{array}{c|c}
k&a_k(t)\\ \hline
+1&H^{-1}\cosh(Ht)\\
0&a_\ast e^{Ht}\\
-1&H^{-1}\sinh(Ht)
\end{array}
\]
en sus dominios respectivos.

\begin{proposition}[Una expansión, tres curvaturas espaciales]
\label{prop:tres-foliaciones-rev6}
Las tres funciones anteriores satisfacen
\[
 \frac{\dot a_k^2+k}{a_k^2}=H^2,\qquad
 \frac{\ddot a_k}{a_k}=H^2,
\]
y por tanto
\[
 \boxed{\quad
 R^{(4)}=12H^2,\qquad
 R^{(3)}=\frac{6k}{a_k^2}.
 \quad}
\]
\end{proposition}

\begin{proof}
Para \(k=+1\) se usa
\(\tanh^2+\operatorname{sech}^2=1\); para \(k=-1\),
\(\coth^2-\operatorname{csch}^2=1\); la carta plana es inmediata.
Las fórmulas de curvatura de FLRW dan entonces los dos escalares.
\end{proof}

El signo \(k\) distingue la geometría de las hojas espaciales; no cambia la
curvatura del espacio--tiempo de de Sitter.  Este ejemplo impide confundir
la lectura de una hoja con la del objeto total, el mismo error que la
arquitectura APP--TRIT--TPK evita desde el comienzo.

\section{\(2\) representaciones operatorias exactas y \(1\) realización geométrica global de la memoria de frontera}

\begin{theorem}[Una memoria, dos representaciones y una lectura global]
\label{thm:una-memoria-tres-funciones}
La separación par--impar del transporte de frontera admite dos
representaciones operatorias exactas y una lectura geométrica global:
\[
\begin{array}{c|c|c|c}
\text{lector}&\text{sector impar}&\text{sector par}&\text{estatuto}\\ \hline
\text{interior}&\text{circulación}&\text{contracción/coercividad}
&\text{representación exacta}\\
\text{reflejado}&\text{holonomía}&\text{kernel positivo/gap finito}
&\text{representación exacta}\\
\text{global}&\text{orientación de la hoja}&
                \text{respuesta media a la expansión}
&\text{lectura FLRW tipada}
\end{array}
\]
Las dos primeras filas comparten la identidad de Stokes y el espectro no
nulo de \(d^\ast d\) y \(dd^\ast\).  La tercera conserva la distinción entre
curvatura de las hojas y curvatura global, pero no define \(\zeta\) desde el
operador dodecafásico. Ninguna fila se obtiene renombrando los símbolos de
otra.
\end{theorem}

\begin{proof}
La primera y la segunda filas son los
\cref{thm:contraccion-flujo-finito,thm:coercividad-gap-gauge-finito,thm:hoja-comun-frontera}.
La estructura local/no local se prueba en la
\cref{prop:green-local-memoria-rev6}.  La tercera fila se realiza por la
descomposición FLRW y la \cref{prop:tres-foliaciones-rev6}; estas fórmulas
prueban la separación geométrica, no un selector autónomo de la viscosidad
de bulto. En cada caso, la involución de orientación fija el canal impar y
la comparación con la ruta inversa fija el canal par.
\end{proof}

La conclusión no es que un fluido, una conexión y una cosmología sean el
mismo sistema físico.  Es más precisa: las tres ciencias reciben
representaciones distintas de un único sistema operatorio de memoria.  El flujo la
lee hacia el interior, la teoría gauge hacia la frontera y la cosmología a
escala global.
